% !TITLE 浣溪沙·谁念西风独自凉
% !AUTHOR 纳兰性德
% !DYNASTY 明清
% !GENRE 词
% !ORDER 4

\begin{poem}
谁念西风\textsuperscript{1}独自凉，萧萧\textsuperscript{2}黄叶闭疏窗\textsuperscript{3}，沉思往事立残阳。

被酒\textsuperscript{4}莫惊春睡重，赌书消得泼茶香\textsuperscript{5}，当时只道是寻常\textsuperscript{6}。
\end{poem}

\begin{annotations}
\notetext{1}{西风：秋风。西风起，天气转凉，独处的凄凉感油然而生。}
\notetext{2}{萧萧：风吹树叶的沙沙声，也形容萧索寂寥。}
\notetext{3}{疏窗：窗格稀疏的窗户。关上窗户也挡不住秋天的凉意。}
\notetext{4}{被酒：酒醉。“被”通“披”，沾染之意。酒醉后不要惊醒沉沉的春睡。}
\notetext{5}{赌书消得泼茶香：用李清照《金石录后序》典故。李清照与丈夫赵明诚常比赛记书，指某事在某书某卷某行，说对的人先喝茶。赢的人往往笑得把茶泼在衣服上——“赌书消得泼茶香”。消得，值得、消受得起。}
\notetext{6}{当时只道是寻常：当时只觉得是稀松平常的小事，如今回想才知那些日子多么珍贵。这是全词最痛的一句，也是纳兰词中最广为流传的名句。}
\end{annotations}

\begin{appreciation}
这是纳兰性德悼念亡妻卢氏的词，中国悼亡词里最平常、也最扎心的一首之一。

"谁念西风独自凉"——秋风起了，还有谁会想着我冷不冷？一个反问，没有答案：没有人了。卢氏嫁给他三年就去世了，从此再没有人问他添衣。这个"凉"字，一半是秋风，一半是心寒。

"萧萧黄叶闭疏窗"，落叶沙沙地响，他关上窗。关上窗挡不住秋天，也挡不住回忆。"沉思往事立残阳"——一个人站在夕阳里想往事，九个字一幅画，画里没有人陪他。

下阕写两件往事。"被酒莫惊春睡重"——有一回妻子喝醉了，春睡沉沉，他轻手轻脚，不敢惊动。"赌书消得泼茶香"用的是李清照的典故：李清照和丈夫赵明诚比赛背书，说某句在某人某书的第几卷第几页第几行，说中的先喝茶，赢的人往往笑得把茶泼了一身。书里读过李清照的词，她后来也写了悼亡——世上恩爱的夫妻都是同一个模子：先是泼茶的笑，后是独自的凉。

最后一句是全词最痛的地方："当时只道是寻常。"那时候只觉得是平平常常的日子。他不写妻子多好，不写自己多想，只说了这一句——当时没觉得稀罕，如今回不去了。苏轼写"十年生死两茫茫"，是放声大哭；纳兰不哭，只说一句"寻常"，声音越轻，越扎人。

妈妈做饭、你写作业的晚上，也是这样的"寻常"日子。这样的晚上，过了就没有了，你多记着点——趁还来得及，别等将来只能说一句"当时只道是寻常"。
\end{appreciation}
